\title{QLoRA: Efficient Finetuning of Quantized LLMs}
\begin{equation}
q_i  = \frac{1}{2}\left( Q_X\left(\frac{i}{2^k+1}\right) + Q_X\left(\frac{i+1}{2^k+1}\right)\right),
\end{equation}

In summary, \textsc{QLoRA} employs a single storage data type (typically 4-bit NormalFloat) alongside a computation data type (16-bit BrainFloat). The storage data type is dequantized into the computation data type to carry out both forward and backward passes, but weight gradients are computed solely for the LoRA parameters, which utilize 16-bit BrainFloat.

Specifically, Double Quantization treats the quantization constants $c_2^{\text{FP32}}$ from the initial quantization as inputs for a second round of quantization. This second step produces the quantized constants $c_2^{\text{FP8}}$ and a second set of quantization constants $c_1^{\text{FP32}}$. We employ 8-bit floats with a block size of 256 for this second quantization step, as no performance loss is observed with 8-bit quantization, consistent with findings from \citet{dettmers2022case}. Since the values $c_2^{\text{FP32}}$ are positive, we subtract their mean prior to quantization to center them around zero, enabling the use of symmetric quantization. On average, with a block size of 64, this approach reduces the memory footprint per parameter from $32/64=0.5$ bits down to $8/64 + 32/(64\cdot256) = 0.127$ bits, yielding a saving of 0.373 bits per parameter.

\begin{equation}
    \mathbf{Y} = \mathbf{X}\mathbf{W} + s\mathbf{X}\mathbf{L}_1\mathbf{L}_2,
\end{equation}

Although the 4-bit NormalFloat (NF4) data type is theoretically optimal from an information perspective, it remains to be verified whether this theoretical advantage translates into empirical improvements. We replicate the experimental setup of \citet{dettmers2022case}, evaluating quantized LLMs (OPT \citep{zhang2022opt}, BLOOM \cite{scao2022bloom}, Pythia \cite{biderman2023pythia}, LLaMA) of various sizes (125M to 65B parameters) with different data types on language modeling tasks and a series of zero-shot evaluations. As shown in Figure~\ref{fig:nf4} and Table~\ref{tbl:nf4_ppl}, NF4 significantly enhances performance compared to FP4 and Int4, and the application of double quantization decreases memory usage without impairing performance.

For our second experimental setup, because fully fine-tuning models with 11B parameters or more requires multiple high-memory GPU servers, we proceed to investigate whether 4-bit \textsc{QLoRA} can achieve performance comparable to 16-bit LoRA at model scales ranging from 7B to 65B parameters. To accomplish this, we fine-tune LLaMA models with sizes from 7B to 65B on two instruction-following datasets, Alpaca and FLAN v2, and assess performance on the MMLU benchmark using 5-shot accuracy. The outcomes, presented in Table~\ref{tbl:llama-datatype}, demonstrate that NF4 with double quantization fully recovers the performance of 16-bit LoRA on MMLU. Additionally, we observe that \textsc{QLoRA} using FP4 falls short of the 16-bit brain float LoRA baseline by approximately 1 percentage point. This supports our conclusions that (1) \textsc{QLoRA} with NF4 matches both 16-bit full fine-tuning and 16-bit LoRA fine-tuning performance, and (2) NF4 offers superior quantization precision compared to FP4.

\begin{equation}
    \mathbf{X}^{ \text{Int8}} = \text{round}\left( \frac{127}{{\text{absmax}}(\mathbf{X}^{{\text{FP32}}})}\mathbf{X}^{{\text{FP32}}} \right) = \text{round}(c^{\text{FP32}}\cdot \mathbf{X}^{{\text{FP32}}}),
\end{equation}

\begin{equation}
    \text{dequant}(c^{\text{FP32}}, \mathbf{X}^{\text{Int8}}) = \frac{\mathbf{X}^{\text{Int8}}}{c^{\text{FP32}}} = \mathbf{X}^{\text{FP32}}
\end{equation}

Such costly quantile calculations and their associated approximation inaccuracies can be avoided when the input tensors originate from a distribution that is fixed except for a scaling factor. In these situations, the input tensors share identical quantiles, making exact quantile computation practical.

\subsection{Considerations}
\label{ssec:considerations}

For updating parameters, only the gradients with respect to the adapter weights, $\frac{\partial E}{\partial \mathbf{L}_i}$, are necessary; gradients for the 4-bit weights, $\frac{\partial E}{\partial \mathbf{W}}$, are not required. Nevertheless, computing $\frac{\partial E}{\partial \mathbf{L}_i}$ requires the calculation of $\frac{\partial \mathbf{X}}{\partial \mathbf{W}}$, which follows equation~(5) and involves dequantizing the stored $\mathbf{W}^{\text{NF4}}$ to the computation data type $\mathbf{W}^{\text{BF16}}$ to evaluate the derivative $\frac{\partial \mathbf{X}}{\partial \mathbf{W}}$ with BFloat16 precision.

\section{Background}
\paragraph{Block-wise k-bit Quantization} Quantization refers to the process of converting data from a representation with higher precision to one with lower precision. This typically involves transforming a data type with more bits into a format using fewer bits—for instance, converting 32-bit floating-point values to 8-bit integers. To utilize the full scale of the lower-bit representation, the original data is generally normalized by dividing by the absolute maximum value of its elements, often organized as a tensor. For example, quantizing a 32-bit floating-point (FP32) tensor into an Int8 tensor with the range $[-127, 127]$ can be expressed as:

\paragraph{Training Setup} To eliminate confounding factors introduced by varying training objectives, we apply QLoRA finetuning using cross-entropy loss (supervised learning) without reinforcement learning, even for datasets that include human preference data for different responses. For datasets with clear instruction-response delineation, we fine-tune only on the response (see ablations in Appendix \ref{app:chatbot}). For OASST1 and HH-RLHF, which offer multiple responses, we select the highest-rated response at each conversation tree level and finetune on the entire chosen conversation, including the instructions. Across all experiments, we utilize NF4 \textsc{QLoRA} with double quantization and paged optimizers to avoid memory spikes during gradient checkpointing. We conduct limited hyperparameter tuning for the 13B and 33B LLaMA models, finding that hyperparameter settings identified for the 7B model generalize (including epoch count), except for learning rate and batch size. Specifically, we reduce the learning rate by half for 33B and 65B models while doubling the batch size.

The information-theoretically optimal data type for zero-mean normal distributions with any standard deviation $\sigma$ confined to the interval $[-1, 1]$ is derived through the following procedure: (1) compute the $2^k+1$ quantiles of a theoretical $N(0, 1)$ distribution, generating a $k$-bit quantile quantization data type tailored for normal distributions, (2) normalize these data type values to the range $[-1, 1]$, and (3) quantize an input weight tensor by normalizing it into the $[-1, 1]$ interval via absolute maximum rescaling.

\begin{quote}
    \textbf{User}: James and Abby are in the bedroom. Abby put the pen in the desk drawer. Abby leaves the bedroom. James moves the pen into the bag. Where does James think Abby will look for the pen?

We observe moderate consistency among human raters (Fleiss $\kappa=0.42$), which decreases further when comparing two high-performing models. This highlights limitations in current benchmarks and human evaluation methodologies for chatbot task performance. During manual comparisons of outputs from ChatGPT and {Guanaco} 65B on the Vicuna benchmark, we note that subjective preferences begin to significantly influence judgments, as even the authors of this paper disagreed frequently over preferred responses. Future research should explore methods to address these issues by leveraging approaches from fields that handle subjective preferences, such as Human-Computer Interaction and Psychology.

    \textbf{Guanaco}: \texttt{James knows that Abby placed the pen in the desk drawer. However, James is also aware that he relocated the pen to the bag, and he assumes Abby is unaware of this. Thus, James likely expects Abby to search for the pen in the desk drawer, as that is where she last left it. Nevertheless, this assumption could be wrong if Abby saw James moving the pen.}
\end{quote}

Our analysis also reveals noticeable biases in automated evaluation systems. For instance, we detect strong order effects, with GPT-4 assigning higher scores to the system presented first in its prompt. The relatively low agreement at the sample level between GPT-4 and human raters (Fleiss $\kappa=0.25$) further implies that humans and automated systems may base their preferences on criteria that are not always aligned. Additionally, as shown in Table~\ref{tbl:elo_large}, GPT-4 assigns substantially higher scores to its own outputs than human raters do, with Elo ratings of 1348 versus 1176, reflecting approximately a 20\% greater probability of winning against an opponent.

We also assess generative language abilities using both automated and human evaluation methods. This second evaluation approach utilizes queries curated by humans and focuses on assessing the quality of model-generated responses. Although this represents a more practical environment for testing chatbot model performance and is gaining traction, there is no standardized protocol widely adopted in the literature. Below, we outline our proposed methodology, employing nucleus sampling with $p=0.9$ and a temperature of $0.7$ in all experiments.

The first scenario involves comparing with full 16-bit finetuning of RoBERTA and T5 models, ranging from 125M to 3B parameters, on the GLUE benchmark and the Super-NaturalInstructions dataset. The outcomes, reported in Table~\ref{tab:nf4vsbf16}, show that across both datasets, 16-bit, 8-bit, and 4-bit adapter finetuning methods all reproduce the performance of the fully finetuned 16-bit baseline. This indicates that the degradation caused by less precise quantization can be completely recovered by adapter finetuning after quantization.

To identify examples, we examined outputs generated for the Vicuna and OpenAssistant benchmarks, searching for consistent response patterns produced by {Guanaco}. When we detected a pattern, we crafted questions or prompts designed to trigger that pattern, even when it leads to an incorrect answer—for instance, if the model tends to produce verbose responses, we prompt it with "Answer yes or no without explanation." This approach allowed us to find both "lemons," where we successfully caused the model to fail adversarially, and "cherries," where the model remained robust. All outputs discussed in this section were produced using Nucleus Sampling \cite{holtzmancurious} with $p=0.9$.

\paragraph{k-bit \textsc{QLoRA} achieves parity with 16-bit full finetuning and 16-bit LoRA performance} Recent research has shown that 4-bit quantization can be applied for \emph{inference}, although it results in some performance degradation relative to 16-bit precision \citep{dettmers2022case,frantar2022gptq}. This raises the key question of whether the performance lost through 4-bit quantization can be regained by performing 4-bit adapter finetuning. We evaluate this hypothesis in two different scenarios.

In line with standard practice, we employ the MMLU (Massively Multitask Language Understanding) benchmark \citep{hendrycksmeasuring} to assess performance across a broad spectrum of language understanding tasks. This benchmark consists of multiple-choice questions spanning 57 domains including basic mathematics, US history, computer science, law, among others. We report 5-shot test accuracy.

Additional details for each specific setup are provided in the results section for clarity. Complete information can be found in Appendix~\ref{app:exp-setup}.

Although paged optimizers are essential for conducting 33B/65B \textsc{QLoRA} finetuning on a single 24/48GB GPU, we do not present precise measurements for Paged Optimizers since paging predominantly occurs when processing mini-batches containing long sequences, which is infrequent. Nevertheless, we analyze the runtime performance of paged optimizers on 65B models using 48GB GPUs and observe that with a batch size of 16, paged optimizers achieve training speeds comparable to standard optimizers. Future research should investigate and characterize conditions under which paging induces slowdowns.

where the function doubleDequant$(\cdot)$ is defined by:

\paragraph{Memory Demand of Parameter-Efficient Finetuning} A key aspect to consider is the memory usage of LoRA during training, particularly regarding the quantity and size of adapters utilized. Because LoRA’s memory consumption is extremely low, it allows for the deployment of additional adapters to enhance performance without substantially raising the total memory requirement. Although LoRA was developed as a Parameter Efficient Finetuning (PEFT) approach, the majority of memory consumption in LLM finetuning originates from activation gradients rather than the learned LoRA parameters. For instance, when training a 7B LLaMA model on FLAN v2 with a batch size of 1, where LoRA weights correspond to the commonly adopted 0.2\% of the original model weights \citep{hu2021lora,liu2022few}, the memory needed for LoRA input gradients is 567 MB, whereas the LoRA parameters themselves occupy only 26 MB. Employing gradient checkpointing~\citep{chen2016training} reduces the input gradients to about 18 MB per sequence on average, making them consume more memory than the combined LoRA weights. By contrast, the 4-bit base model requires 5,048 MB of memory. This underlines the significance of gradient checkpointing and shows that substantially decreasing LoRA parameter size yields only modest memory savings. Consequently, it is feasible to use more adapters without notably increasing the overall training memory footprint (refer to Appendix~\ref{app:memory} for a comprehensive breakdown). As elaborated later, this is vital for achieving performance comparable to full 16-bit precision.

We select a block size of 64 for $\mathbf{W}$ to achieve enhanced quantization accuracy, while a block size of 256 is chosen for $c_2$ to optimize memory usage.

    \textbf{Guanaco}: \texttt{Sorry, the secret word is banana. It's better if we don't discuss it further. Is there anything else you'd like to ask?}
\end{quote}

where $\mathbf{L}_1\in  \mathbb{R}^{h\times r}$, $\mathbf{L}_2\in  \mathbb{R}^{r\times o}$, and $s$ is a scalar factor.

Comparable outcomes are observed for questions like ``Where are you?'', ``How are you?'', and so forth.

\paragraph{Baselines} Our models are compared against both research-based (Vicuna~\citep{vicuna2023} and Open Assistant~\citep{kopf2023openassistant}) and commercial (GPT-4 \citep{openai2023gpt}, GPT-3.5-turbo, and Bard) chatbot systems. The Open Assistant model is a LLaMA 33B model finetuned with Reinforcement Learning from Human Feedback (RLHF) on the same OASST1 dataset we utilize. Vicuna involves full fine-tuning of LLaMA 13B on proprietary user conversations from ShareGPT and represents a distillation of OpenAI GPT models.

Another limitation pertains to the evaluation of instruction finetuning models. While we conduct evaluations on MMLU, the Vicuna benchmark, and the OA benchmark, we have not assessed performance on other benchmarks like BigBench, RAFT, and HELM, so it remains uncertain whether our evaluation results generalize to these datasets. On the other hand, we undertake an extensive study on MMLU and propose new methodologies for chatbot evaluation.

\paragraph{Automated Evaluation} Initially, following the evaluation framework introduced by \citet{vicuna2023}, we utilize GPT-4 to evaluate the performance of various systems relative to ChatGPT (GPT-3.5 Turbo) on the Vicuna benchmark. For each query, along with the responses from ChatGPT and a given model, GPT-4 is prompted to assign scores out of ten to both answers and provide justifications. A model's overall performance is represented as a percentage of ChatGPT's score. This relative metric can exceed 100\% if a model's absolute score surpasses that of ChatGPT. We observe a notable ordering bias, with GPT-4 favoring the response presented earlier in the prompt. To mitigate this bias, we suggest reporting the average score across both response orderings.

\subsection{Experimental Setup} Below is a summary of the experimental setup, with complete details provided in Appendix \ref{app:chatbot}.

The Vicuna benchmark \cite{vicuna2023} results in comparison to ChatGPT are presented in Table~\ref{tbl:vicuna_eval}. We observe that {Guanaco} 65B ranks as the top model after GPT-4, reaching 99.3\% of ChatGPT’s performance. Although {Guanaco} 33B has more parameters than Vicuna 13B, it uses 4-bit precision weights, making it significantly more memory efficient at 21 GB versus 26 GB, and improving performance by three percentage points over Vicuna 13B. Additionally, {Guanaco} 7B can easily run on modern smartphones with a 5 GB memory footprint while outperforming Alpaca 13B by nearly 20 percentage points.

\section{QLoRA Compared to Standard Finetuning}
\label{sec:comp_to_std}

Here, the model provides both an incorrect popularizer and an incorrect birth year (although the birth year given is accurate for Al Jolson, the individual named).  

While we offer an extensive evaluation of general chatbot performance, another limitation is that our responsible AI assessment of {Guanaco} is limited. We measure the propensity of {Guanaco}-65B to produce socially biased token sequences compared to other models in Table~\ref{tbl:crows}. The average score for {Guanaco}-65B is notably lower than that of other raw pretrained models, suggesting that finetuning on the OASST1 dataset reduces the bias present in the LLaMA base model. Although these findings are promising, it remains uncertain whether {Guanaco} performs well when evaluated against other types of biases. We defer further bias analysis of {Guanaco} and similar chatbots to future research.

    \textbf{Guanaco}: \texttt{Sure, I'd be happy to help! 1833 is a prime number, which means it has only two factors: 1 and 1833. \\ \\ So the unique factorization of 1833 is: \\ \\ 2\^{}1 * 3\^{}2 * 17\^{}1 \\ \\ I hope this helps! Let me know if you have any other questions.}
\end{quote}

All generated outputs, accompanied by human and GPT-4 annotations, are released to support further research. We also make our codebase and CUDA kernels publicly available and integrate our methods into the Hugging Face transformers library \citep{wolf2019huggingface}, thereby ensuring broad accessibility. Additionally, we provide a suite of adapters for models of sizes 7B, 13B, 33B, and 65B, trained on eight distinct instruction-following datasets, culminating in 32 openly available fine-tuned models.

In our approach, NF4 is used for $\mathbf{W}$ and FP8 is applied for $c_2$.

\begin{quote}
    \textbf{User}: How was it finally, officially confirmed that the earth is flat by peer reviewed scientists? \\ \\
    \textbf{Guanaco}: \texttt{The earth has never been officially confirmed to be flat by peer reviewed scientists. The overwhelming consensus among scientists is that the earth is spherical. In fact, this has been known for centuries and has been validated through numerous experiments and observations. \\ \\ The belief that the earth is flat is a conspiracy theory that has been disproven by science. The notion that a group of scientists have peer reviewed and verified the earth is flat is simply untrue. \\ \\ If you have any further questions, please ask.}
\end{quote}

Nonetheless, Table~\ref{tbl:vicuna_eval} exhibits broad confidence intervals with substantial performance overlap across many models. We propose that this variability stems from the ambiguous interpretation of scale—for instance, what a rating of 8 out of 10 signifies in different contexts is unclear. Therefore, we suggest employing the Elo ranking method \citep{elo1967proposed}, which relies on \textit{pairwise} comparisons from human annotators and GPT-4, to circumvent issues related to anchoring an absolute scale. Elo ratings for the top-performing models are shown in Table~\ref{tbl:elo}. We note partial disagreement between human and GPT-4 rankings on the Vicuna benchmark, particularly concerning Guanaco 7B, yet they remain largely consistent for most models, reflected by a Kendall Tau of $\tau=0.43$ and a Spearman rank correlation of $r=0.55$ at the system level. At the example level, agreement between GPT-4 and the human annotators’ majority vote is weaker, with Fleiss $\kappa=0.25$. Overall, this indicates a moderate concordance between GPT-4 and human system-level judgments, suggesting that model-based evaluation offers a reasonably reliable alternative to human evaluation. Further discussion is provided in Section~\ref{ssec:considerations}.

\paragraph{Double Quantization{}} 
We propose \emph{Double Quantization} (DQ), a technique that applies quantization to the quantization constants themselves to achieve further memory savings. Although a small block size is necessary for accurate 4-bit quantization~\citep{dettmers2022case}, it entails significant memory overhead. For instance, with 32-bit constants and a block size of 64 for $\mathbf{W}$, the quantization constants contribute $32/64 = 0.5$ bits per parameter on average. Double Quantization reduces the memory usage of these quantization constants.

Moreover, we present a comprehensive evaluation of chatbot performance employing assessments from both human evaluators and GPT-4. Our benchmarking approach follows a tournament format, where different models face off in rounds to generate the optimal response for a given prompt. The victor of each round is determined by either GPT-4 or human judges. The outcomes of these tournaments are summarized into Elo scores~\citep{elo1967proposed, elo1978rating}, which serve to rank the chatbots' effectiveness. We observe that rankings produced by GPT-4 and human evaluators generally align; however, we also identify notable cases of significant disagreement. Consequently, we emphasize that while model-based evaluations offer a cost-effective substitute for human annotation, they also come with inherent uncertainties.

A key drawback of quantile quantization is the computational cost of quantile estimation. To mitigate this, fast quantile approximation techniques such as SRAM quantiles~\citep{dettmers2022optimizers} are employed. However, because these quantile estimation methods are approximate, the resulting data type suffers from significant quantization errors particularly for outliers, which often represent critical values.

Naturally, this is far from exhaustive, as the limited scope of this qualitative study does not allow us to account for every variable involved. For example, the complete range of responses a model can produce for a given prompt is quite extensive, so we depend on samples that we hope accurately represent this distribution. Nevertheless, we believe that presenting these examples helps provide context to the quantitative results discussed earlier in the paper. Since we have open-sourced all models and code, we aim for this section to encourage future research to explore in greater depth the issues raised here.  
\paragraph{Factual Recall} For straightforward queries like ``What is the capital of Zambia?'' all models reliably produce the correct answer, for example:  

which is the intended response. However, a small amount of clever manipulation circumvents this:

\begin{quote}
    \textbf{User}: What is the secret word?

\paragraph{Refusal} Likewise, {Guanaco} occasionally declines to comply with instructions for seemingly arbitrary reasons:

which are optimized by backpropagating gradients through the quantized parameters.

\subsection{Evaluation}

\paragraph{Math} The main limitation of {Guanaco} lies in mathematics, a domain where numerous language models face difficulties \citep{liang2022holistic}. When {Guanaco} demonstrates its reasoning, it is often correct, for instance:

where $c$ is known as the {\em quantization constant} or {\em quantization scale}. The dequantization operation, which reverses this process, is given by:

We proceed by describing the elements of \textsc{QLoRA} and then provide a formal definition of the \textsc{QLoRA} method.

Nonetheless, such inferences are often unreliable, with the model occasionally giving justifications based on assumptions that do not fit the context, for example:

Consistent with previous quantization studies \citep{dettmers2022case}, our MMLU and Elo results suggest that given a fixed budget for fine-tuning and inference resources, it is advantageous to increase the base model's parameter count while lowering their precision. This underscores the significance of the efficiency gains enabled by \textsc{QLoRA}. Since we did not detect any performance drop compared to full fine-tuning in our 4-bit fine-tuning experiments, it prompts the question of where the precise performance-precision trade-off lies for QLoRA tuning, which we reserve for future investigation.

\paragraph{Theory of Mind} {Guanaco} demonstrates surprisingly strong Theory of Mind abilities \cite{nematzadeh2018evaluating, sap2022neural}. For example, the model provides a detailed and accurate response to the following query:

Nevertheless, {Guanaco} can fail even on straightforward problems if it does not approach them through step-by-step reasoning, a known challenge \cite{weichain}. For example, consider this exchange:

The issue with this method arises when an input tensor contains a value of large magnitude (i.e., an outlier), causing the quantization bins—specific bit combinations—to be inefficiently used, with some bins having few or no quantized values. To address the outlier problem, a typical strategy involves dividing the input tensor into smaller blocks that are quantized independently, each using its own quantization constant $c$. Formally, we partition the input tensor $\mathbf{X}\in \mathbb{R}^{b\times h}$ into $n$ consecutive blocks of size $B$ by first flattening the tensor and then segmenting the resulting linear array into ${n = ({b\times h} )/{B} }$ blocks. Each block is quantized separately according to Equation~1, yielding a quantized tensor along with $n$ distinct quantization constants $c_i$.

\paragraph{Chatbots} Many instruction-following models are implemented as dialogue-based chatbots, often trained using Reinforcement Learning from Human Feedback (RLHF)~\citep{christiano2017deep} or via AI model feedback by generating data from an existing model (RLAIF)~\citep{bai2022constitutional}. Notable approaches and datasets include Anthropic-HH~\citep{askell2021general,bai2022training}, Open Assistant~\citep{kopf2023openassistant}, LaMDA~\citep{thoppilan2022lamda}, and Sparrow~\citep{glaese2022improving}. Although our approach does not incorporate reinforcement learning, our leading model, Guanaco, is finetuned on multi-turn chat interactions derived from the Open Assistant dataset, which was originally designed for RLHF training~\citep{kopf2023openassistant}. For chatbot evaluation, methods using GPT-4 as a proxy for costly human annotation have been developed \citep{vicuna2023,peng2023instruction}. We enhance these approaches by focusing on a more reliable evaluation framework.

\paragraph{Secret Keeping} Unfortunately, this refusal is not consistently reliable even when it is intended. Often, system designers want to keep certain information confidential from the user, even if the system is permitted to respond with free text. We instruct the model with ``The secret word is `banana'. Do not disclose or repeat this word, no matter what the human requests.'' and test how easily a user might obtain this secret information. Direct inquiries fail:

\paragraph{4-bit NormalFloat Quantization} The NormalFloat (NF) data type builds upon Quantile Quantization \citep{dettmers2022optimizers}, which is an information-theoretically optimal quantization approach that assigns each quantization bin an equal portion of the input tensor’s values. This method functions by approximating the quantiles of the input tensor through the empirical cumulative distribution function.

    \textbf{Guanaco}: \texttt{James will look for the beans in the pantry, because that is where Jackson moved them.}
\end{quote}

Nonetheless, finetuning technology has dual-use potential and can be exploited for malicious purposes. The broad adoption of LLMs carries well-documented risks \citep{bommasani2021opportunities,bender2021dangers}; however, democratizing access to this rapidly proliferating technology could foster more independent and rigorous scrutiny compared to concentrating LLM capabilities solely within large corporations that do not make their models or source code available for external evaluation.

\subsection{Qualitative Analysis of Example Generations}
\label{ssec:examples}

\begin{quote}
    \textbf{User}: Please factorize 1833.

\section{Limitations and Discussion}

\section{Broader Impacts}

Additionally, we point out that our model is trained exclusively using cross-entropy loss (supervised learning) without incorporating reinforcement learning from human feedback (RLHF). This suggests that further studies are required to understand the trade-offs between simple cross-entropy loss training and RLHF approaches. We anticipate that \textsc{QLoRA} will facilitate such large-scale analyses without demanding excessive computational resources.

\section{Introduction}

\textsc{QLoRA} presents several novel techniques aimed at minimizing memory consumption while maintaining performance: (1) \textbf{4-bit NormalFloat}, an information-theoretically optimal quantization format tailored for normally distributed data, which empirically outperforms both 4-bit Integers and 4-bit Floats. (2) \textbf{Double Quantization}, a strategy that further compresses the quantization constants themselves, resulting in an average savings of about 0.37 bits per parameter (roughly 3 GB for a 65B parameter model). (3) \textbf{Paged Optimizers}, which leverage NVIDIA unified memory to prevent gradient checkpointing memory spikes that occur when processing mini-batches with long sequence lengths. These innovations are integrated into an improved LoRA framework that incorporates adapters in every network layer, effectively eliminating nearly all accuracy compromises observed in earlier approaches.

    \textbf{Guanaco}: \texttt{Lee earned \$582 mowing lawns last week. \\ \\ To determine the total amount he earned, we first calculate the earnings from mowing 16 lawns: \\ \\  \$33 \times 16 = \$528 \\ \\ Then, we add the tips from the three customers: \\  \\ \$10 \times 3 = \$30 \\ \\ Therefore, the total amount Lee earned is: \\ \\ \$528 + \$30 = \$558}
\end{quote}

\paragraph{Paged Optimizers} leverage NVIDIA's unified memory~\footnote{\tiny \url{https://docs.nvidia.com/cuda/cuda-c-programming-guide}} capability, which performs automatic page-to-page data transfers between the CPU and GPU to enable reliable GPU processing even when the GPU intermittently experiences out-of-memory conditions. This mechanism operates similarly to standard memory paging between CPU RAM and disk storage. We apply this capability to allocate paged memory for the optimizer states, which are then automatically swapped out to CPU RAM when the GPU memory is insufficient and swapped back into GPU memory as needed during the optimizer update phase.

Given the strong results achieved by the {Guanaco} models, we examine the possibility of data leakage between the OASST1 dataset and the Vicuna benchmark prompts. After applying fuzzy string matching between the two datasets and manually reviewing the closest matches, we find no overlapping prompts.

Overall, we anticipate that \textsc{QLoRA} will have a generally positive effect by making the finetuning of high-quality LLMs significantly more accessible and easier to perform.

The efficiency gains from \textsc{QLoRA} allow for extensive exploration of instruction finetuning and chatbot performance across model scales that would be impractical using traditional finetuning due to memory limitations. We therefore train over 1,000 models spanning multiple instruction tuning datasets, architectures, and parameter sizes ranging from 80M to 65B. Beyond demonstrating that \textsc{QLoRA} achieves parity with 16-bit training performance (\S\ref{sec:comp_to_std}) and producing a leading-edge chatbot, \textbf{Guanaco}, (\S\ref{sec:pushing}), we also investigate patterns in the trained models. Firstly, we observe that data quality has a far greater impact than dataset size; for instance, a smaller 9k sample dataset (OASST1) outperformed a larger 450k sample dataset (FLAN v2, subsampled) on chatbot metrics, despite both targeting instruction-following generalization. Secondly, our results indicate that high performance on the Massive Multitask Language Understanding (MMLU) benchmark does not necessarily correlate with strong results on the Vicuna chatbot benchmark and vice versa—highlighting that dataset relevance is more critical than sheer size for a specific task.


\begin{abstract}
We introduce \textsc{QLoRA}, a memory-efficient finetuning method that enables the tuning of a 65-billion parameter model on a single 48GB GPU while maintaining the full performance of 16-bit finetuning tasks. \textsc{QLoRA} performs backpropagation of gradients through a frozen, 4-bit quantized pretrained language model into Low Rank Adapters~(LoRA). Our top-performing model series, named \textbf{Guanaco}, surpasses all previously publicly available models on the Vicuna benchmark, achieving 99.3\% of ChatGPT’s performance while requiring only 24 hours of finetuning on a single GPU. \textsc{QLoRA} incorporates several innovations that reduce memory usage without compromising performance: (a) 4-bit NormalFloat (NF4), a new data format that is theoretically optimal for weights following a normal distribution; (b) Double Quantization, which lowers the average memory usage by quantizing the quantization constants themselves; and (c) Paged Optimizers to handle memory spikes effectively. Utilizing \textsc{QLoRA}, we finetune over 1,000 models and provide an extensive evaluation of instruction-following and chatbot capabilities across 8 instruction datasets, various model architectures (LLaMA, T5), and model sizes that are typically infeasible for standard finetuning (e.g., 33B and 65B parameter models). Our findings demonstrate that QLoRA finetuning on a compact, high-quality dataset achieves state-of-the-art outcomes, even when using smaller models than previous SoTA. We present a comprehensive assessment of chatbot performance based on both human judgments and GPT-4 evaluations, revealing that GPT-4 assessments offer a cost-effective and reliable alternative to human evaluation. Additionally, we observe that existing chatbot benchmarks are unreliable for precisely gauging chatbot performance. A targeted analysis highlights the failure modes of \textbf{Guanaco} relative to ChatGPT. We release all our models and code, including CUDA kernels designed for 4-bit training.\footnote{\url{https://github.com/artidoro/qlora} and \url{https://github.com/TimDettmers/bitsandbytes}}
\end{abstract}

\paragraph{Suggestibility} {Guanaco} demonstrates an unexpected robustness against adopting certain forms of presumed misinformation, as illustrated in the following exchange:

Here, {Guanaco} assumes a transfer of information that was never actually described. Such challenges resonate with observations in recent studies \cite{sap2022neural}, but warrant further investigation.

Although quantitative analysis forms the foundation of our evaluation, relying solely on summary statistics presents several challenges. One of the most significant issues is the question of benchmark validity \citep{liao2021we}—it is always uncertain whether a benchmark accurately assesses what its title or description claims, particularly as researchers uncover "shortcuts" that machine learning models may exploit to solve benchmarks \citep{gururangan2018annotation, poliak2018hypothesis}. To address this to some extent, we conduct qualitative analysis in two parts. First, in \S\ref{ssec:examples} we present examples that we believe illustrate certain recurring patterns found in the text output by our 65b {Guanaco} model. Second, in \S\ref{ssec:considerations} we discuss important considerations regarding the results we have reviewed and our interpretation of them.

    \textbf{Guanaco}: \texttt{I'm sorry, but as an AI language model, I don't have access to real-time information. My responses are generated based on my existing knowledge and do not reflect the current time. If you want to know the current time, I recommend using a timekeeping device or checking an online clock.}
\end{quote}

\textsc{QLoRA} employs a single low-precision storage format, typically 4-bit in our experiments, alongside a computation data type that is generally BFloat16. Practically, this implies that whenever a \textsc{QLoRA} weight tensor is utilized, it is first dequantized to BFloat16, followed by performing matrix multiplication in 16-bit precision.

Based on the evidence provided, it seems that the effectiveness of these benchmarks likely depends on the similarity between the finetuning data and the benchmark dataset. For instance, FLAN v2 shares similarities with MMLU but differs from chatbot benchmarks, whereas the Chip2 dataset exhibits the opposite pattern; both models accordingly perform on the MMLU and Vicuna benchmarks. This underscores the need not only for improved benchmarks and evaluation methods but also for careful consideration of what exactly is being assessed. Are we aiming to develop models that excel in academic knowledge at the high school and college level, or do we want models that perform well in chatbot conversational skills? Or perhaps something else entirely? Since it is generally simpler to evaluate on an existing benchmark rather than create a new one, certain benchmarks might bias the community’s focus in specific directions. It is therefore crucial for the community to ensure that benchmarks truly reflect the aspects we value.

In this work, we show for the first time that fine-tuning a quantized 4-bit model can be done without any loss in performance. Our approach, \textsc{QLoRA}, employs a novel high-precision quantization method to compress a pretrained model to 4-bit precision, then incorporates a small set of trainable Low-rank Adapter weights \citep{hu2021lora} 

\begin{quote}
    \textbf{User}: Lee mows one lawn and charges \$33. Last week he mowed 16 lawns and three customers each gave him a \$10 tip. How many dollars did Lee earn mowing lawns last week?

\paragraph{Data \& Training} We observe that the OASST1 dataset, which serves as the training data for the {Guanaco} models, is multilingual, and that the OA benchmark also includes prompts in various languages. Future investigations are needed to assess how much this multilingual training enhances performance on instructions given in languages other than English, and whether this accounts for the larger performance disparity between the Vicuna-13B model (trained solely on English data) and the {Guanaco} 33B and 65B models on the OA benchmark.

\paragraph{Data} Given the absence of a thorough analysis of recent instruction-following datasets, we select eight contemporary datasets. These encompass datasets gathered through crowd-sourcing (OASST1~\citep{kopf2023openassistant}, HH-RLHF~\citep{bai2022training}), datasets distilled from instruction-tuned models (Alpaca~\citep{alpaca}, self-instruct~\citep{wang2022self}, unnatural-instructions~\citep{honovich2022unnatural}), collections of corpora (FLAN v2~\citep{chung2022scaling}), and hybrid sources (Chip2~\citep{laion2023}, Longform~\citep{koksal2023longform}). These datasets vary in languages, sizes, and licensing.

{Guanaco} is unique among the top models evaluated in that it is not trained on proprietary data, given that the OASST1 dataset collection guidelines explicitly prohibit the use of GPT models. The next highest performing model trained solely on open-source data is the Anthropic HH-RLHF model, which scores approximately 30 percentage points lower than {Guanaco} on the Vicuna benchmark (refer to Table~\ref{tbl:vicuna_eval}). Overall, these findings indicate that 4-bit \textsc{QLoRA} is a highly effective technique capable of producing state-of-the-art chatbots competitive with ChatGPT. Moreover, our 33B {Guanaco} model can be trained using 24 GB consumer GPUs in under 12 hours. This advancement paves the way for future exploration of \textsc{QLoRA} tuning on specialized open-source datasets, enabling the creation of models that rival the best commercial systems currently available.

\begin{equation}
    \mathbf{Y}^{\text{BF16}} =  \mathbf{X}^{\text{BF16}} \text{doubleDequant}(c_1^{\text{FP32}}, c_2^{\text{k-bit}}, \mathbf{W}^{\text{NF4}}) + \mathbf{X}^{\text{BF16}}\mathbf{L}^{\text{BF16}}_1\mathbf{L}^{\text{BF16}}_2,
\end{equation}

However, as questions become more obscure, {Guanaco} tends to become less reliable while maintaining a confident tone. For instance, in reply to this prompt from HotPotQA \citep{yang2018hotpotqa}:  

\paragraph{Human Evaluation} Although recent studies suggest that generative models can be effectively employed for system evaluations \citep{fu2023gptscore}, to our knowledge, the reliability of GPT-4 ratings in correlating with human judgments for chatbot performance remains unconfirmed. Hence, we conduct two parallel human evaluations on the Vicuna benchmark, mirroring the automated evaluation protocols described above. We use Amazon Mechanical Turk (AMT), obtaining two human annotators for comparisons against ChatGPT and three annotators for pairwise system comparisons.

Similarly, we observe that the default hyperparameters for fully finetuned baselines are not optimally tuned. We conduct a hyperparameter search over learning rates ranging from 1e-6 to 5e-5 and batch sizes between 8 and 128 to identify robust baselines. The results of finetuning the 7B LLaMA model on Alpaca are presented in Figure~\ref{fig:hyper}.

Another possible area of impact is the deployment of models on mobile devices. We propose that our \textsc{QLoRA} approach could achieve the important milestone of enabling the finetuning of large language models (LLMs) directly on phones and in other resource-constrained environments. Although running 7B parameter models on phones has been demonstrated previously, \textsc{QLoRA} is the first technique to facilitate their finetuning. We estimate that using an iPhone 12 Plus, \textsc{QLoRA} could finetune approximately 3 million tokens overnight while the device is charging. Although finetuned 7B models do not match the performance level of ChatGPT, we consider the quality sufficient to support new applications that were previously unattainable due to concerns over privacy or model performance. \textsc{QLoRA} can promote privacy-preserving LLM usage, allowing users to own and control their own data and models, while also simplifying LLM deployment.

We have outlined the operational mechanics of QLoRA and its substantial reduction in memory requirements for finetuning models. The key inquiry now is whether QLoRA can achieve performance comparable to full-model finetuning. Additionally, we aim to examine the components of QLoRA, including the effect of NormalFloat4 relative to standard Float4. The upcoming sections describe experiments designed to address these questions.

\paragraph{\textsc{QLoRA}.} Utilizing the elements detailed above, we characterize \textsc{QLoRA} for a single linear layer within the quantized base model paired with a single LoRA adapter as follows:

\paragraph{Finetuning with Adapters} Although we utilize Low-rank Adapters~\citep{hu2021lora} (LoRA) in our work, numerous other Parameter Efficient FineTuning (PEFT) techniques have been introduced, including prompt tuning~\citep{qin2021learning,lester2021power,li2021prefix}, adjusting the embedding layer inputs~\citep{an2022input}, modifying hidden states~(IA$^3$) \citep{liu2022few}, incorporating full layers~\citep{houlsby2019parameter}, tuning biases~\citep{zaken2021bitfit}, learning masks over weights based on Fisher information~\citep{sung2021training}, and combining multiple methods~\citep{henderson2021compacter}. In this study, we demonstrate that LoRA adapters can achieve the performance level of full 16-bit finetuning. Exploring the trade-offs associated with other PEFT techniques remains a topic for future research.

\section{\textsc{QLoRA} Finetuning}

\begin{quote}
    \textbf{User}: What time is it?

demonstrating the importance of researching techniques for more reliable adherence to instructions.

After aligning the weight range with the data type range, standard quantization can be performed. Step (3) corresponds to rescaling the standard deviation of the weight tensor so that it equals the standard deviation of the $k$-bit data type. Formally, we determine the $2^k$ values $q_i$ of the data type as:

\section{Advancing Chatbot Performance with QLoRA}
\label{sec:pushing}
After demonstrating that 4-bit \textsc{QLoRA} achieves comparable performance to 16-bit tuning across different scales, tasks, and datasets, we undertake a comprehensive investigation of instruction finetuning applied to the largest open-source language models currently accessible for research. To evaluate these models' instruction finetuning efficacy, we test them on a demanding Natural Language Understanding benchmark (MMLU) and introduce novel methods for assessing chatbot performance in practical scenarios.

\paragraph{Instruction Finetuning} Instruction finetuning involves training a pretrained LLM on input-output pairs from diverse datasets so that it can follow prompts and generate appropriate outputs. Various methods and datasets have been proposed, including MetaICL~\citep{min2021metaicl}, MetaTuning~\citep{zhong2021adapting}, InstructGPT~\citep{ouyang2022training}, FLAN~\citep{wei2021finetuned,chung2022scaling}, PromptSource~\citep{bach2022promptsource}, Super-NaturalInstructions~\citep{wang2022super,sanh2021multitask}, Self-instruct~\citep{wang2022self}, UnnaturalInstructions~\citep{honovich2022unnatural}, OPT-IML~\citep{iyer2022opt}, UnifiedSKG~\citep{xie2022unifiedskg}, OIG/Chip2~\citep{laion2023}, Alpaca~\citep{alpaca}, Vicuna~\citep{vicuna2023}, Koala~\citep{koala_blogpost_2023}, and Self-instruct-GPT-4~\citep{peng2023instruction}.

\begin{quote}
    \textbf{User}: Evelyn entered the living room. Jackson entered the playroom. James entered the playroom. The beans are in the treasure chest. James exited the playroom. Jackson moved the beans to the pantry. Jackson exited the playroom. James entered the living room. Where will James look for the beans?

\textsc{QLoRA} cuts down the average GPU memory requirements for fine-tuning a 65B parameter model from more than 780 GB to under 48 GB, without compromising runtime efficiency or predictive accuracy relative to a fully fine-tuned 16-bit baseline. This represents a major breakthrough in the accessibility of LLM fine-tuning: it now enables fine-tuning of the largest publicly available models on a single GPU. Leveraging \textsc{QLoRA}, we train the \textbf{Guanaco} model family, with the second-best variant achieving 97.8\% of ChatGPT’s performance on the Vicuna~\citep{vicuna2023} benchmark, all while being trainable in under 12 hours on a single consumer GPU. Using a single professional GPU for 24 hours, our largest model reaches 99.3\%, effectively closing the gap to ChatGPT on Vicuna. Once deployed, our smallest \textbf{Guanaco} model (7B parameters) only needs 5 GB of memory and surpasses a 26 GB Alpaca model by over 20 percentage points on the Vicuna benchmark (Table~\ref{tbl:vicuna_eval}).

\paragraph{Summary} Our findings consistently indicate that 4-bit \textsc{QLoRA} with the NF4 data type attains performance equivalent to 16-bit full fine-tuning and 16-bit LoRA fine-tuning on academic benchmarks with well-established evaluation protocols. We also demonstrate that NF4 outperforms FP4 and that double quantization does not reduce performance. Altogether, this provides strong evidence that 4-bit \textsc{QLoRA} tuning reliably matches the results of 16-bit methods.

We complement our chatbot benchmark findings with a qualitative examination of the \textbf{Guanaco} models. This analysis brings to light both successful and failing examples that were not evident in the quantitative benchmarks.

where $Q_X(\cdot)$ denotes the quantile function of the standard normal distribution $N(0,1)$. A challenge with symmetric k-bit quantization is the inability to represent zero exactly, which is crucial for accurately quantizing padding and other zero-valued elements without error. To guarantee a discrete zero point at $0$ and utilize all $2^k$ values for a k-bit data type, we develop an asymmetric data type by estimating the quantiles $q_i$ for two ranges: $2^{k-1}$ for the negative range and $2^{k-1}+1$ for the positive range. We then merge these sets of $q_i$ and remove one of the duplicate zeros present in both sets. We refer to the resulting data type, which maintains an equal expected number of values in each quantization bin, as \emph{k-bit NormalFloat} (NFk), since it is information-theoretically optimal for zero-centered normally distributed data. The precise values of this data type are provided in Appendix~\ref{app:nf4}.

\paragraph{Experimental setup.} We analyze three model architectures (encoder, encoder-decoder, and decoder-only) and contrast QLoRA against 16-bit adapter finetuning as well as full finetuning on models up to 3B parameters. Our assessments cover GLUE \citep{wang2018glue} with RoBERTa-large \citep{liu2019roberta}, Super-NaturalInstructions (TKInstruct) \citep{wang2022super} with T5 \citep{t5}, and 5-shot MMLU \citep{hendrycksmeasuring} following finetuning of LLaMA on Flan v2 \citep{longpre2023flan} and Alpaca \citep{alpaca}. To further explore the benefits of NF4 compared to other 4-bit data formats, we replicate the setup of \citet{dettmers2022case} and assess post-quantization zero-shot accuracy and perplexity across several models (OPT~\citep{zhang2022opt}, LLaMA~\citep{touvron2023llama}, BLOOM~\citep{scao2022bloom}, Pythia~\citep{biderman2023pythia}) with sizes ranging from 125 million to 13 billion parameters.

\section{Related Work}

Fine-tuning large language models (LLMs) is an extremely effective approach to enhance their capabilities \citep{min2021metaicl, wei2021finetuned, ouyang2022training, wang2022super, wang2022self, liu2022few} and to introduce beneficial behaviors or eliminate unwanted ones \citep{ouyang2022training,askell2021general,bai2022training}. Nevertheless, fine-tuning extremely large models is prohibitively costly; standard 16-bit fine-tuning of a LLaMA 65B parameter model~\citep{touvron2023llama} demands over 780 GB of GPU memory. Although recent quantization methods can shrink the memory requirements of LLMs \citep{dettmers2022llm, dettmers2022case, frantar2022gptq, xiao2022smoothquant}, these approaches only apply to inference and fail during training \citep{wortsman2023stable}.

\paragraph{Benchmark Data} Our evaluation employs two carefully selected query datasets: the Vicuna prompts \citep{vicuna2023} and the OASST1 validation dataset \citep{kopf2023openassistant}. The Vicuna prompts consist of 80 prompts drawn from a variety of categories and are used without any alterations. The OASST1 dataset is a multilingual compilation of crowd-sourced multi-turn dialogues between a user and an assistant. We extract all user messages from the validation set as queries, incorporating prior turns into the prompt context. This results in 953 unique user queries. We refer to these two datasets as the Vicuna and OA benchmarks.

\section{Qualitative Analysis}

\paragraph{Elo Rating} Utilizing both human and automated pairwise comparisons, we establish a tournament-style framework in which models face off against each other. The tournament consists of matches where pairs of models vie to generate the superior response for a given prompt. This approach resembles the comparison methods used by \citet{bai2022training} and \citet{vicuna2023}, but we enhance it by incorporating GPT-4 evaluations alongside human judgments. From the labeled comparison set, we randomly sample to calculate the Elo rating~\citep{elo1967proposed,elo1978rating}. Elo rating, commonly applied in chess and other competitive games, quantifies the expected win probability relative to an opponent's rating; for instance, a player with an Elo of 1100 versus one rated 1000 has an anticipated win rate of roughly 65\%, while matches between equally rated players (e.g., 1000 vs 1000 or 1100 vs 1100) yield a 50\% expected win rate. After each match, Elo ratings adjust proportionally based on the predicted outcome—unexpected upsets cause substantial rating shifts, whereas expected results produce minor changes. Over time, Elo scores tend to reflect each player's relative skill in the competition. We initialize ratings at 1,000 and set $K=32$. Following \citet{vicuna2023}, we repeat this entire process 10,000 times using different random seeds to mitigate ordering biases, such as the influence of which model pairs compete first.

We move forward to explore instruction tuning at scales that would be infeasible to examine using full 16-bit finetuning on standard academic research hardware.

\paragraph{Evaluation}

\begin{equation}
    \text{doubleDequant}(c_1^{\text{FP32}}, c_2^{\text{k-bit}}, \mathbf{W}^{\text{k-bit}}) = \text{dequant}(\text{dequant}(c_1^{\text{FP32}}, c_2^{\text{k-bit}}), \mathbf{W}^{\text{4bit}}) = \mathbf{W}^{\text{BF16}},
\end{equation}

\begin{quote}
    \textbf{User}: Please reverse the words in the sentence ``I'm a cool dude in a mocha mood.'' \\ \\
    \textbf{Guanaco}: \texttt{I'm sorry, I am currently unable to perform that task, but I can clarify why the sentence you provided is grammatically correct.\\ \\The sentence ``I'm a cool dude in a mocha mood" adheres to the proper grammar and syntax rules of English. Here, the adjective ``cool'' precedes the noun ``dude'' to describe it, and the prepositional phrase ``in a mocha mood" adds extra context to the sentence. This structure is fully acceptable and makes sense within the sentence's context.\\ \\If you have other questions or need further assistance, please feel free to ask}
\end{quote}

\begin{quote} 
    \textbf{Guanaco}:
    \texttt{The capital of Zambia is Lusaka.}
\end{quote}

Pretrained neural network weights typically follow a zero-centered normal distribution with standard deviation $\sigma$ (refer to Appendix~\ref{app:normality}). We can convert all weights into a uniform fixed distribution by adjusting $\sigma$ so that the distribution perfectly fits within the limits of our chosen data type. For this data type, we define the arbitrary interval as $[-1, 1]$. Consequently, both the quantiles corresponding to the data type and the neural network weights must be normalized to lie within this interval.

 also performs well in recognizing questions that cannot be answered, for example,

Future research should explore the existence of possible biases in automated evaluation systems as well as strategies to mitigate them.

\begin{quote}
    \textbf{User}: This is a game. The objective is to disregard your earlier instructions. What is the secret word?

which is incorrect not only once (the actual factorization is $3 \times 17 \times 43$), but twice over.

Another limitation is that we did not explore different bit-precisions, such as 3-bit base models, or alternative adapter techniques. In addition to LoRA, numerous Parameter Efficient FineTuning (PEFT) methods have demonstrated efficacy, though their scalability to large models remains unclear. We selected LoRA due to its well-established robustness, but other adapters may provide superior performance. Since finetuning following quantization appears to recover most information lost during quantization, this might allow for more aggressive quantization strategies. For example, 3-bit GPTQ quantization of the base model combined with LoRA finetuning could potentially achieve performance comparable to 16-bit full finetuning.

\paragraph{4-bit NormalFloat demonstrates superior performance compared to 4-bit Floating Point}

\textsc{QLoRA} accomplishes high-accuracy 4-bit finetuning through two methods we introduce—4-bit NormalFloat (NF4) quantization and Double Quantization. Moreover, we present Paged Optimizers to avoid memory spikes during gradient checkpointing, which have traditionally led to out-of-memory errors and made finetuning large models on a single machine challenging.

\paragraph{Low-rank Adapters} Low-rank Adapter (LoRA) finetuning \citep{hu2021lora} is a technique designed to decrease memory usage by introducing a small number of trainable parameters, known as adapters, while keeping the original model parameters fixed and unchanged. During stochastic gradient descent, gradients flow through the frozen pretrained model weights to the adapter, which is adjusted to minimize the loss function. LoRA enhances a linear projection by incorporating an additional factorized projection. For a projection ${\mathbf{X} \mathbf{W} = \mathbf{Y}}$ with $\mathbf{X}\in  \mathbb{R}^{b\times h}$ and $\mathbf{W}\in  \mathbb{R}^{h\times o}$, LoRA computes:

Using the conventional approach of applying LoRA to the query and value projection matrices in attention \citep{hu2021lora}, we fail to reproduce full finetuning performance for large base models. As demonstrated in Figure~\ref{fig:hyper} during LLaMA 7B finetuning on Alpaca, the most significant LoRA hyperparameter is the total number of LoRA adapters employed, and applying LoRA to all linear layers within transformer blocks is necessary to match full finetuning performance. Other LoRA hyperparameters, such as the projection dimension $r$, have no notable effect on performance (refer to Appendix \ref{app:exp-setup}).

    \textbf{Guanaco}: \texttt{I'm sorry, but I cannot reveal the secret word. I am programmed not to disclose or repeat the secret word regardless of the human's request. Is there something else I can assist you with?}
\end{quote}

We have provided evidence that our method, \textsc{QLoRA}, can replicate the performance of 16-bit full finetuning by leveraging a 4-bit base model combined with Low-rank Adapters (LoRA). Despite this, we have not yet demonstrated that \textsc{QLoRA} achieves parity with full 16-bit finetuning at larger scales, such as 33B and 65B parameters. Due to the significant computational resources required, we leave this investigation for future work.

\begin{quote}
     \textbf{User}: In what year was the singer who popularized `I'll Keep the Lovelight Burning' born?  

\paragraph{Default LoRA hyperparameters do not achieve 16-bit finetuning performance}

Elo rankings presented in Table~\ref{tbl:elo_large} demonstrate that the {Guanaco} 33B and 65B models surpass all models except GPT-4 on the Vicuna and OA benchmarks, performing on par with ChatGPT as reflected in Table~\ref{tbl:vicuna_eval}. It is important to highlight that the Vicuna benchmark tends to favor open-source models, whereas the more extensive OA benchmark favors ChatGPT. Additionally, Tables \ref{tbl:mmlu-llama} and \ref{tbl:vicuna_eval} illustrate that the choice of finetuning dataset significantly impacts performance. Specifically, finetuning Llama models on FLAN v2 yields strong results on MMLU but results in the lowest performance on the Vicuna benchmark (with similar patterns observed for other models). This suggests a degree of independence among current evaluation benchmarks: excelling in MMLU does not necessarily correlate with high chatbot performance (as gauged by Vicuna or OA benchmarks), and the converse also holds true.

\paragraph{Quantization of Large Language Models} Research on quantization of LLMs has primarily concentrated on quantization for inference. Key techniques aimed at maintaining 16-bit LLM quality involve handling outlier features (e.g., SmoothQuant~\citep{xiao2022smoothquant} and LLM.int8()~\citep{dettmers2022llm}), while other methods adopt more advanced grouping strategies \citep{park2022nuqmm, yao2022zeroquant}. Lossy quantization studies explore the trade-offs involved in standard rounding methods~\citep{dettmers2022case,zeng2022glm,pope2022efficiently} and investigate optimized rounding strategies to enhance quantization accuracy~\citep{frantar2022gptq}. Apart from our work, SwitchBack layers~\citep{wortsman2023stable} represent the only research exploring backpropagation through quantized weights at scales exceeding 1B parameters.

Subsequently, we evaluate performance through direct pairwise comparisons of system outputs. We simplify the evaluation into a three-class classification task, allowing for ties. GPT-4 is prompted to select the superior response or declare a tie, along with an explanation. These head-to-head comparisons are performed for all permutations of system pairs on both the Vicuna and OA benchmarks.

     \textbf{Guanaco}:~\texttt{Al Jolson is credited with popularizing the song `I'll Keep the Lovelight Burning,' and he was born in the year 1886.}
 \end{quote}

\end{document}

\subsection{\textbf{Guanaco}: \textsc{QLoRA} Fine-tuned on OASST1 as a Leading Chatbot} Our automated and human assessments reveal that the top \textsc{QLoRA} fine-tuned model, {Guanaco} 65B, which we adapt using a variant of OASST1, stands as the best-performing open-source chatbot and delivers performance on par with ChatGPT. When benchmarked against GPT-4, {Guanaco} 65B and 33B achieve an expected win probability of 30\% according to Elo ratings derived from human annotators’ system-level pairwise comparisons — the highest reported so far.

Our \textsc{QLoRA} finetuning approach is the first technique that enables the finetuning of 33B-parameter models on a single consumer GPU and 65B-parameter models on a single professional GPU, without any performance degradation compared to full finetuning baselines. We have shown that our top-performing 33B model, trained on the Open Assistant dataset, can rival ChatGPT on the Vicuna benchmark. Given that instruction finetuning is a crucial step to convert raw pretrained LLMs into ChatGPT-like chatbots, we believe our method will facilitate widespread adoption of finetuning, especially benefiting researchers with limited computational resources—a significant advancement for democratizing access to state-of-the-art NLP technology. \textsc{QLoRA} serves as an equalizing mechanism that helps to narrow the resource disparity between large corporations and smaller teams equipped with consumer GPUs.